\documentclass[11pt]{article}
\usepackage{amsmath,amssymb,amsthm}
\newtheorem{proposition}{Proposition}
\newtheorem{lemma}[proposition]{Lemma}
\title{Auditing the Pair-Odd Coefficient Compiler}
\author{}
\date{}
\begin{document}
\maketitle

\paragraph{Scope.}
The pair-odd kernel is a valid new coefficient-domain construction.
The results here verify its bounded expander span and identify a concrete
failure of the proposed small-tail argument for transferring coefficient
signs to signs of cube values. None of these statements is Lean-checked.
No attention definition or admissibility condition is changed.

\section*{The exact two-head family}
Take three coordinate pairs $(x,x'),(y,y'),(z,z')$, a real $R>1$, and
\[
 A=x-x',\quad B=y-y',\quad
 D=1+R(2-x-x')+R(2-y-y')+(2-z-z').
\]
For $0<t<1$ define
\[
 F_t=\frac1{2t}\left(\frac{A}{D-tB}-\frac{A}{D+tB}\right)
     =\frac{AB}{D^2-t^2B^2}.
 \tag{1}
\]
All denominator slopes are strictly negative, and the denominators are
positive throughout the Boolean cube: $D\ge1$ and $|tB|<1$.
Each summand is therefore an original admissible affine-ratio head by
\texttt{FracAtom.ofNegativeAffineRatio}. On this finite cube,
$F_t\to F_0=AB/D^2$ uniformly as $t\to0$.

Swapping $x,x'$ negates $F_t$, and swapping $y,y'$ also negates it.
Uniqueness of the multilinear representative implies that any coefficient
whose support contains both or neither members of the row pair is zero;
the same statement holds for the column pair. In particular, all
coefficients of degrees six and five vanish exactly.

\begin{proposition}[Unbounded lower tails inside the admissible family]
Let $K_4,K_3,K_2$ be the coefficients of $xyz z'$, $xyz$, and $xy$
in the multilinear representative of $F_0$. Then
\[
\begin{aligned}
 K_4&=\frac1{(2R+1)^2}-\frac2{(2R+2)^2}+\frac1{(2R+3)^2},\\
 K_3&=\frac1{(2R+2)^2}-\frac1{(2R+3)^2},\\
 K_2&=\frac1{(2R+3)^2}.
\end{aligned}
 \tag{2}
\]
All three are positive, and
\[
 K_4\sim\frac3{8R^4},\qquad K_3\sim\frac1{4R^3},\qquad
 K_2\sim\frac1{4R^2}.
 \tag{3}
\]
Consequently $K_3/K_4\sim2R/3$ and $K_2/K_4\sim2R^2/3$.
These unbounded ratios also occur in actual two-head scores $F_t$,
with all the stated exact upper-coefficient zeros.
\end{proposition}
\begin{proof}
For $D>0$, $D^{-2}=\int_0^\infty u e^{-Du}\,du$.
Multilinear coefficient extraction is a finite linear combination of
Boolean evaluations, so it commutes with this integral. For example,
the $xyz z'$ coefficient is
\[
 \int_0^\infty u e^{-(2R+1)u}(1-e^{-u})^2\,du.
\]
Deleting one or both reserve coordinates gives respectively
\[
 \int_0^\infty u e^{-(2R+2)u}(1-e^{-u})\,du,
 \qquad \int_0^\infty u e^{-(2R+3)u}\,du.
\]
Integration yields (2), and expansion at $R=\infty$ yields (3).
For each fixed $R$, coefficient extraction is continuous in the finite
cube values. Thus the same ratios are approached by $F_t$ as $t\to0$.
Choose $R$ first and a sufficiently small positive $t$ second.
\end{proof}

This disproves a uniform claim that such modules can be normalized to
have unit desired coefficient and arbitrarily small lower coefficients.
It does not disprove a different compiler that cancels or controls the
lower tails by a new mechanism.

\section*{Why the coefficient-sign bridge is not yet a proof}
If a multilinear polynomial is
$P(x)=\sum_T c_Tx^T$, its value at $\mathbf1_S$ is
$\sum_{T\subseteq S}c_T$. Therefore coefficient signs alone do not
determine cube-value signs. For instance $P(x)=2-x$ has negative
$x$ coefficient but a positive value at the vertex $x=1$.

One sufficient condition is
$|c_S|>\sum_{T\subsetneq S}|c_T|$ for every nonempty $S$, with the
desired coefficient signs. Descending-degree synthesis does not
automatically produce this condition. Exact zero coefficients above
the programmed degree preserve coefficients set earlier, but earlier
modules already create lower coefficients. To change a later
coefficient's sign, its new contribution may have to cancel a large
existing value. It cannot simply be assigned an arbitrarily small
amplitude. Equations (2)--(3) exhibit this issue inside the proposed
two-head family.

Nor can one obtain dominance by substituting a large number for each
Boolean variable in the rational formula. Boolean multilinearization
and dilation do not commute. Already
\[
 F(x)=\frac1{1+x},\quad
 P(x)=1-\frac{x}{2}\quad(x\in\{0,1\}),\qquad
 F(3)=\frac14>0>P(3)=-\frac12.
 \tag{4}
\]
A valid transfer would have to control the actual cube values, or
implement the dilated multilinear polynomial by admissible heads with
the asserted cost. Equation (4) invalidates substituting directly into
the rational heads as an implementation of that dilation.

\section*{A bounded expander span does hold}
The separate incidence-matrix issue has a positive answer. Let $U$ and
$V$ be oriented vertex-edge incidence matrices of finite graphs on
$r$ and $s$ vertices. Let $P_r=I-\mathbf1\mathbf1^\top/r$ and define
$P_s$ similarly. Assume the scalar $\ell^1$ Poincare inequalities
\[
 \|P_r a\|_1\le C_r\|U^\top a\|_1,
 \qquad
 \|P_s b\|_1\le C_s\|V^\top b\|_1.
 \tag{5}
\]

\begin{proposition}[Uniformly bounded edge-edge coefficients]
If an $r\times s$ matrix $Y$ has every row sum and every column sum
zero, then there is a matrix $K$, indexed by the two edge sets, with
\[
 Y=UKV^\top,\qquad
 \|K\|_\infty\le C_rC_s\|Y\|_\infty.
 \tag{6}
\]
If both graph degrees are at most $d$, entrywise rounding $K$ with
error at most $\delta$ changes $UKV^\top$ entrywise by at most
$d^2\delta$.
\end{proposition}
\begin{proof}
Apply (5) to the columns, and then to the transposes of the rows, to get
for every real $r\times s$ matrix $Z$
\[
 \|P_r ZP_s\|_1
 \le C_r\|U^\top ZP_s\|_1
 \le C_rC_s\|U^\top ZV\|_1.
 \tag{7}
\]
Here matrix $\ell^1$ is the sum of the absolute values of all entries.
Since $P_rYP_s=Y$, the Frobenius pairing obeys
\[
 |\langle Y,Z\rangle|
 =|\langle Y,P_rZP_s\rangle|
 \le C_rC_s\|Y\|_\infty\|U^\top ZV\|_1.
 \tag{8}
\]
Define $\ell(U^\top ZV)=\langle Y,Z\rangle$ on the image of the
linear map $Z\mapsto U^\top ZV$. Equation (8) proves both that this
definition is independent of the choice of $Z$ and that its norm
for the entrywise $\ell^1$ norm is at most
$C_rC_s\|Y\|_\infty$. Finite-dimensional Hahn--Banach extends it
to all edge-edge matrices with the same bound. The dual of entrywise
$\ell^1$ is entrywise $\ell^\infty$, so the extension is pairing
with a matrix $K$ of the asserted norm. For all $Z$,
\[
 \langle Y,Z\rangle=\langle K,U^\top ZV\rangle
                  =\langle UKV^\top,Z\rangle,
\]
which proves (6). Finally each output entry sums at most $d^2$
signed entries of $K$, proving the rounding assertion.
\end{proof}

For completeness, an edge-isoperimetric lower bound
$|\partial S|\ge h|S|$ for $|S|\le r/2$ implies (5) with
$C_r=2/h$. Choose a median $m$ of $a$. The layer-cake formula and
the edge bound give
$\sum_i|a_i-m|\le h^{-1}\sum_{ij\in E}|a_i-a_j|$.
Since
$r|\overline a-m|\le\sum_i|a_i-m|$, the triangle inequality gives
$\sum_i|a_i-\overline a|\le2\sum_i|a_i-m|$.
Thus bounded-degree expanders give constants independent of the number
of vertices. Two perfect matchings whose union is a cycle give the
algebraic zero-sum span, but not these uniform constants.

The span theorem addresses only the double-centered coefficient
matrix. Mean and row/column components, and the coefficient-to-value
bridge discussed above, still require their own valid constructions.
\end{document}
